% !TeX root = ../main.tex
\section{Conclusion}
\label{sec:conclusion}

We presented a behavioral simulation framework for analog compute-in-memory inference that disentangles algorithmic generalization failures from physical precision-retention constraints. By comparing ideal pure-quantization baselines against realistic phase-change memory (PCM) physics, the framework serves as a practical decision aid for establishing deployment frontiers in low-precision analog hardware.

Several observations stand out. At extreme low precision (4-bit), fixed hardware-instance transfer is the most prominent simulated deployment risk observed in our matrix: naive analog deployment collapses to 14.64\%, whereas Ensemble Hardware-Aware Training (HAT) resamples mismatch masks to rescue convergence, recovering fresh-instance accuracy to 86.16$\pm$0.19\%. Under the tested PCM UnitCell simulation regime, convergence is observed across the 4/6/8-bit precision ladder where pure quantization fails. However, this convergence exposes a precision-retention tradeoff. While 8-bit PCM is drift-flat and 4-bit PCM reaches the most aggressive tested compression at the cost of measurable time-dependent drift, 6-bit emerges as the best observed Pareto midpoint in this sweep, maintaining 8-bit-like stability alongside robust fresh accuracy.

Ultimately, this framework clarifies that standard fixed-mask training is insufficient for extreme low-precision analog deployment. Ensemble HAT resolves this algorithmic bottleneck, allowing system designers to navigate the subsequent physical limits of the analog precision-retention frontier. Extending these hardware-aware training principles to memory-bound inference components, such as analog KV-cache, presents a compelling future trajectory for large-scale deployment.

\section*{Data Availability}
All datasets used in this study (CIFAR-10, CIFAR-100, and Flowers-102) are publicly available through standard PyTorch torchvision and HuggingFace dataset interfaces. The source data underlying the main-text and Supplementary figures and tables are packaged in the submission-matched source-data archive together with the corresponding manifests, logs, and figure-generation inputs. A reproducibility archive containing the code snapshot, source data, SHA-256 manifest, and \texttt{CITATION.cff} file has been staged for Zenodo deposition and will be assigned a DOI at acceptance.

\section*{Code Availability}
The simulation framework, training and evaluation scripts, device-profile utilities, and plotting code will be provided to editors and reviewers through a reviewer-accessible archive at submission. A curated public release will be made available at \url{https://github.com/Leslie360/HAT} upon acceptance. The current repository is distributed under the Apache-2.0 licence, and the staged Zenodo archive will be updated with the final public DOI at acceptance.

\section*{Competing Interests}
The authors declare no competing interests.
